% 本文件由 LaTeX 排版 Agent 根据有效 translation.json 自动生成。
% author=Councils; work_id=3811; title=卡尔西顿大公会议（主历451年）

\chapter{卡尔西顿大公会议（主历451年）}

% structure: p0001 source=h1 target=omitted reason=page-title-already-rendered-by-parent

\Emph{（《会议记录》摘录）}\par

% structure: p0003 source=h2 target=section reason=relative-heading-rank; base=h2

\section{第一会期}

最可敬的帕斯卡西努斯主教，宗座代牧，与他的同僚们一同起立，站在会众中间说：我们从至福的、宗徒的罗马城主教——此城是所有教会的元首——处接到指示，指示说，狄奥斯库鲁斯不得在此集会中就座，倘若他试图就座，应被驱逐。我们必须执行此指示；若如今诸位圣座下令，就请他离开，否则我们便离开。\par

最荣耀的审判官和全体元老院说：你对最可敬的狄奥斯库鲁斯主教提出什么特别的指控？\par

最可敬的帕斯卡西努斯主教，宗座代牧，说：既然他已到来，必须对他提出异议。\par

最荣耀的审判官和全体元老院说：按照刚才所说的，请具体说明他所犯的指控。\par

最可敬的卢森提乌斯主教，代表宗座，说：请他为自己的判决说明理由。因为他竟然对他无权管辖的人作出判决。而且他胆敢在未经宗座授权的情况下召开主教会议，这是从未发生过、也不可能发生的事。\par

最可敬的帕斯卡西努斯主教，代表宗座，说：我们不能违背至福的、宗徒的主教（拉丁文本中\Quote{主教}对应\Quote{教宗}）的谕令，他治理着宗座，也不能违背教会法典或教父传统。\par

最荣耀的审判官和全体元老院说：你们应当具体说明他在哪些方面犯了错误。\par

可敬的卢森提乌斯主教，代表宗座，说：我们不会容忍如此大的错误对我们和你们造成伤害，以至于这个前来接受审判的人却就座（如同要作出判决的人一样）。\par

荣耀的审判官和全体元老院说：如果你担任法官的职务，你就不应当像被审判的人那样为自己辩护。\par

在最荣耀的审判官和神圣会议（\OriginalTerm{神圣会议}{\GreekText{τῆς} \GreekText{ἱερᾶς} \GreekText{συγκλήτου}}）的命令下，最热心的亚历山大主教狄奥斯库鲁斯坐在会众中间，最可敬的罗马主教们也在各自的位置坐下，保持沉默。最可敬的多里莱乌姆城主教优西比乌斯走到中间说：\par

［\Emph{随后他提交了一份请愿书，并宣读了\Quote{强盗会议}的记录。同时也宣读了君士坦丁堡会议在弗拉维亚诺主持下针对优提克斯的记录}（第175栏）。］\par

宣读完毕后，最荣耀的审判官和庞大的会议（\OriginalTerm{庞大的会议}{\GreekText{ὑπερφυὴς} \GreekText{σύγκλητος}}）说：在座的最可敬的主教们，你们认为神圣的弗拉维亚诺主教在阐明信仰时，是否保持了正统和公教信仰，还是在这方面有所偏离？\par

最可敬的帕斯卡西努斯主教，代表宗座，说：蒙福的弗拉维亚诺极其神圣而完美地阐明了信仰。他的信仰和阐述与至福的、宗徒的罗马主教的书信一致。\par

最可敬的君士坦丁堡总主教亚纳多留说：蒙福的弗拉维亚诺优美而正统地阐述了我们教父的信仰。\par

最可敬的卢森提乌斯主教，宗座代牧，说：既然蒙福的弗拉维亚诺的信仰与宗座和教父传统一致，那么公正的做法是，他被异端者定罪的判决，应由这至圣会议转回他们身上。\par

最可敬的叙利亚安提约基雅主教玛克西穆说：蒙福的弗拉维亚诺总主教正统地阐明了信仰，并符合最属神的、至圣的莱奥总主教。我们所有人都热切地接受这一点。\par

最可敬的卡帕多细亚凯撒勒雅主教塔拉修斯说：蒙福的弗拉维亚诺是按照蒙福的济利禄所说的而发言的。\par

［\Emph{于是，主教们一个接一个地表达了他们的意见。随后继续宣读君士坦丁堡会议记录。}］\par

在宣读的这一点上，最可敬的亚历山大总主教狄奥斯库鲁斯说：我接受\Quote{出自二性}，但我不接受\Quote{二性}（\OriginalTerm{我接受\Quote{出自二性}，但我不接受\Quote{二性}}{\GreekText{τὸ} \GreekText{ἐκ} \GreekText{δύο} \GreekText{δέχομαι}, \GreekText{τὸ} \GreekText{δύο} \GreekText{οὐ} \GreekText{δέχομαι}}，第371栏）。\par

% structure: p0023 source=h2 target=section reason=relative-heading-rank; base=h2

\section{第三会期}

［\Emph{皇室代表似乎没有出席，在君士坦丁堡的副主教埃提乌斯宣布开会之后，}］\par

来自西西里省利利贝乌姆的帕斯卡西努斯主教，代表至圣的旧罗马宗座总主教莱奥，用拉丁语发言，翻译如下：本属神的会议深知，神圣信函已被送至至福的、宗徒的教宗莱奥，邀请他屈尊莅临圣会议。但由于古老习俗不允许，且当时的普遍必要性似乎也不许可，我的卑微代表他\OriginalTerm{托付于圣会议的事}{\GreekText{τὰ} \GreekText{τῆς} \GreekText{ἁγίας} \GreekText{συνόδου} \GreekText{ἐπέτρεψε}}，因此必须通过我的\OriginalTerm{干预}{\GreekText{διαλαλιᾶς}}来审查所有提出讨论的事项。［拉丁文本在我放置普通文本希腊文的地方读作\Quote{命令我的卑微代替他主持本次神圣会议}。］因此，请接受我们最属神的兄弟、同为主教的优西比乌斯所呈上的书卷，并由属神的副主教、首席公证员埃提乌斯宣读。\par

副主教、首席公证员埃提乌斯接过书卷，宣读如下。\par

[接下来是欧瑟伯的请求，以及之后若干内容，四个请求分别致\Quote{至圣、天主所爱的普世总主教及伟大罗马的宗主教良，以及全体在加采东召开的神圣普世主教会议}等。前两个来自亚历山大的执事，第三个来自该教区的一位前司铎，第四个也来自亚历山大的一位平信徒。之后又召集了狄约斯哥若，但他没有来，于是宣判了对他不利的判决，并通过一份列入法令集的信函告知了他。（L. and C., Conc., Tom. IV., col. 418.）主教们大多逐一发表意见，但罗马代表们共同发言，在他们的话中出现如下内容（col. 426:）]\par

因此，至圣和蒙福的良，伟大而更古老的罗马的总主教，通过我们，并经由本次神圣会议，与三倍蒙福、极其荣耀的宗徒伯多禄——他是天主教会磐石和基础，也是正统信仰的基础——一起，剥去了他的主教职位，并使他与一切圣职的尊贵隔绝。因此，愿此至圣而伟大的会议对上述的狄约斯哥若处以教会法律之惩罚。\par

[主教们随后逐一发言支持罢免狄约斯哥若，但通常的理由是他经三次召唤仍拒绝出席。]\par

当所有至圣的主教就此发表意见后，他们签署了以下内容。\par

\Emph{神圣普世会议对狄约斯哥若发出的谴责。}\par

神圣、伟大而普世的会议，因天主的恩宠，按照我们最虔诚和天主所爱的皇帝的命令，在比提尼亚的加采东城，至圣而胜利的殉道者圣欧斐米的殉道堂中召集，致狄约斯哥若。\par

我们谨此告知：在十月十三日，你因藐视神圣教规，不服从本次神圣普世会议，以及你所犯的其他罪行，被神圣普世会议剥夺了主教职位，并被逐出一切教会的\OriginalTerm{秩序}{\GreekText{θεσμοῦ}}，因为即使依照神圣教规，本次神圣大会议三次传唤你回应控告者，你都没有前来。\par

% structure: p0034 source=h2 target=section reason=relative-heading-rank; base=h2

\section{第四次会议}

最显赫而光荣的法官和伟大的元老院说：\par

让可敬的会议现在宣布关于信仰似乎正确的决定，因为已经处理的事已经显明了。帕斯卡西诺和路森提乌斯，最可敬的主教，以及波尼法爵，最可敬的司铎，宗座代表，通过那位最可敬的人，帕斯卡西诺主教说：既然神圣、蒙福和普世的会议持守并追随由尼西亚教父们所确立的\OriginalTerm{信德规则}{fidei regulam}，它也同样确认由可纪念的伟大德奥多西下令召集的150位教父在君士坦丁堡所确立的信仰。此外，可纪念的著名的济利禄在厄弗所会议（其中聂斯多略被定罪）上所发表的信仰阐述也被接受。而第三，那位蒙福之人，所有教会之总主教良的著作显示了真正的信仰，他谴责了聂斯多略和欧迪奇的异端。同样，神圣会议持守这信仰，追随这信仰——它不能增添什么，也不能删减什么。\par

当这些经由神圣御前枢密院的虔诚书记贝罗尼西安译为希腊语后，最可敬的主教们喊道：我们全都相信，我们如此受洗，我们如此施洗，我们如此相信，我们现在如此相信。\par

最光荣的法官和伟大的元老院说：既然我们看到圣福音书被放置在你们的圣座旁，请让这里每一位聚集的主教宣告，至蒙福的良总主教的书信是否与在尼西亚聚集的318位教父的阐述以及后来在皇城聚集的150位教父的法令相符。\par

[对于这个问题，主教们逐一回答，直到收集了161份单独意见，之后帝国法官要求其余主教集体投票（col. 508）。]\par

所有最可敬的主教喊道：我们都同意，我们都如此相信；我们全都同心合意。我们如此心意，我们如此相信，等等。\par

% structure: p0041 source=h2 target=section reason=relative-heading-rank; base=h2

\section{第五次会议}

帕斯卡西诺和路森提乌斯，最可敬的主教，以及司铎波尼法爵，罗马宗座代表说：如果他们不接受那位宗徒性和蒙福之人教宗良的书信，就请下令给我们辞别书，我们好在那里（即西方）召开会议。\par

[随后进行了长时间的辩论，讨论是否应接受拟好并提交的法令。这似乎是大多数主教的想法。最后，委员们提议成立一个由22人组成的委员会，与他们会面并向会议报告，皇帝则以威胁的方式强加此议，否则所有人将被遣送回家，并在西方召集新的会议。即便如此，他们也没有屈服（col. 560）。]\par

最可敬的主教们喊道：皇帝万岁！要么让定义（即本次会议提出的定义）成立，要么我们走。皇帝万岁！\par

塞巴斯托波尔最可敬的主教塞克罗皮乌斯说：我们请求再次宣读定义，凡不赞同且不愿签名的人，可以自行离开；因为我们赞同这些如此优美草拟的内容，并无批评。\par

伊利里亚最蒙福的主教们说：让那些反对的人显明出来。反对的人是聂斯多略派。反对的人，让他们去罗马。\par

最显赫而最光荣的法官说：狄约斯哥若承认他接受了\Quote{两个性体}的说法，但不承认有两个性体。然而至圣的良总主教说，在基督内有两个性体，不变、不可分、无混淆地结合于独生子我们的救主内。你们愿意跟随谁，是至圣的良还是狄约斯哥若？\par

最可敬的主教们喊道：我们相信正如\Person{利奥}。那些矛盾的是欧迪奇派。\Person{利奥}正确地阐述了信德。\par

最尊贵而荣耀的审判官们说：那么，按照我们至圣父亲\Person{利奥}的判断，在定义中添加：在基督内有两个本性，不可改变地、不可分离地、不相混淆地结合。\par

[\Emph{委员会随后在最神圣的殉道者\Person{欧菲米斯}的小堂内开会，之后报告了一篇信德定义，该定义教导同一道理，却不是\WorkTitle{利奥卷}} (col. 562)。]\par

\Emph{加采东大公会议的信德定义。}\par

神圣、伟大、普世的主教会议，藉天主的恩宠和我们的至为虔诚并信奉基督的皇帝\Person{马尔西安}和\Person{瓦伦提尼安}——奥古斯都——的命令，在比提尼亚省都府\Place{加采东}，于神圣而胜利的殉道者\Person{欧菲米亚}的殉道堂内开会，议决如下：\par

我们的主耶稣基督救主，在加强他门徒中信德的认知时，为使任何人都不与邻舍在宗教教义上意见相左，并使真理的宣讲平等地展示给众人，曾说：\Quote{我把平安留给你们，我将我的平安赐给你们。}可是，由于恶者不断在虔敬的种子中撒莠子，并总是发明某种反对真理的新诡计；因此主，如同他始终所做的那样，为人类预备了这一切，兴起了这位虔诚、忠信而热心的君主，并将各地司祭的首领们召集到他身边；这样，在我们共同的主基督的恩宠激励下，我们得以从基督的羊群中驱逐一切谎言的瘟疫，并用真理的柔嫩叶子喂养他们。我们以一致同意做到了这一点，驱除错谬的理论，更新教父们无误的信德，向众人宣告三百一十八位教父的信经，并将接受同一宗教总纲的教父们同等列入其数。这些就是后来在伟大的\Place{君士坦丁堡}聚集、并批准同一信德的一百五十位圣教父。此外，我们遵循圣主教会议在厄弗所先前举行的关于信德的一切规则和形式（由荣福纪念的罗马的\Person{塞来斯定}和亚历山大的\Person{济利禄}领导），我们宣告：由三百一十八位圣而可颂的教父在虔诚纪念的\Person{君士坦丁}治下于\Place{尼西亚}集会上所制定的正确无瑕的信德宣认，应居首位；并且，由一百五十位圣教父在\Place{君士坦丁堡}为根除当时兴起的异端并确认同一公教和宗徒的信德所议决的一切，也应同样有效。\par

\Emph{三百一十八位教父在尼西亚的信经。}\par

我们信唯一的天主，等等。\par

\Emph{同样}，一百五十位圣教父在\Place{君士坦丁堡}聚集的信经。\par

我们信唯一的天主，等等。\par

这一明智而有益的、属神恩宠的公式，已足以完善地认识和确认宗教；因为它教导了关于圣父、圣子、圣神的圆满道理，并藉着忠信地接受它而揭示了主的降生成人。但是，因为那些意图使真理的宣讲落空的人，通过他们各自的异端，引发了空洞的胡言乱语；其中有些人敢于败坏主为我们而降生成人的奥迹，并拒绝使用\OriginalTerm{天主之母}{\GreekText{Θεοτόκος}}来称呼童贞玛利亚，而另一些人则引入混淆和混合，愚昧地认为肉体和神性的本性完全相同，主张独生子的神圣本性通过混合而能够受苦；因此，本次神圣、伟大、普世的主教会议，渴望排除每一个反对真理的诡计，并教导从起初不变的事，一开始就议决：三百一十八位教父的信德应保持完整无缺。并且，为反对那些敌视圣神的人，它确认后来由在帝都聚集的一百五十位圣教父所交付的关于圣神圣体的道理；他们向众人宣告这一道理，并非引入任何先辈所缺乏的东西，而是为反对那些企图摧毁圣神主权的人，通过书面文件解释他们对圣神的信德。而且，为反对那些着手败坏救世工程奥迹并无耻地声称由圣童贞玛利亚所生者只是一个凡人的人，它接受亚历山大教会牧者荣福\Person{济利禄}致\Person{聂斯多略}和东方派的公函，认为这些公函适用于反驳\Person{聂斯多略}的疯狂愚行，并指导那些以神圣热情渴求认识救恩标志的人。此外，为确认正统教义，它正确地在这之上增加了伟大古罗马的主席、至可颂的至圣总主教\Person{利奥}致荣福纪念的\Person{弗拉维安}总主教的信函，以消除\Person{欧迪奇}的错谬教义，认为它符合伟大\Person{伯多禄}的宣认，并且如同反对误信者的共同支柱。因为它反对那些将救世工程奥迹撕裂为两个儿子的人；它从神圣集会中驱逐那些敢于说独生子具有可受苦的神性的人；它抵制那些想象基督两个本性混合或混淆的人；它驱除那些幻想仆人的形状来自天上或某些不是从我们取来的实体的人，并且它诅咒那些愚昧地谈论联合前主有两个本性、并以为联合后只有一个本性的人。\par

我们追随圣教父，异口同声教导人宣认同一圣子［天主］和我们的主耶稣基督为同一个［位格］，他于天主性内圆满，于人性内亦圆满，是真实的天主亦是真实的人，由理性灵魂与［人］身体所构成，就其天主性而论与圣父同体，就其人性而论与我们同体；在一切事上相似我们，惟独没有罪；按其天主性，在万世之前由圣父所生；但在这末世，为了我们人类并为了我们的救恩，由童贞玛利亚——天主之母——按其人性而生于［世界］。这一位同一的耶稣基督、天主的独生子，必须被宣认为具有两种本性，不相混淆、不变换、不分开、不离散［结合］，且本性间的区别绝不因结合而消失，反而每一本性的特性得以保存，并联合于一位格与一自立体中，不分离亦不分裂为两个位格，而是同一圣子、独生子、天主圣言、我们的主耶稣基督，正如古先知关于他所说，并如主耶稣基督所教导我们，也如教父信经所传授给我们。\par

因此，我们以最大的精确与审慎表达上述事项之后，神圣的普世主教会议规定：任何人不得被容许提出另一种信条（\OriginalTerm{另一种信条}{\GreekText{ἑτέραν} \GreekText{πίστιν}}），亦不得撰写、编纂、杜撰或以此教导他人。但若有人胆敢编纂另一种信条，或向那些愿意从外邦人、犹太人或任何异端皈依真理知识的人提出、教导或传授另一种信经（\OriginalTerm{另一种信经}{\GreekText{ἕτερον} \GreekText{σύμβολον}}），若他们是主教或圣职人员，应予以罢免，主教罢免其主教职，圣职人员罢免其圣职；但若他们是隐修士或平信徒，则应处以绝罚。\par

在宣读定义之后，所有最虔诚的主教们喊道：这是教父们的信德；让各大主教立即签署它；请他们立即在法官面前签署它；使这妥善定案之事毫无迟延；这是宗徒们的信德；我们皆以此站立；我们皆如此相信。\par

% structure: p0062 source=h2 target=section reason=relative-heading-rank; base=h2

\section{第六次会议}

\Emph{［皇帝亲临会场并向公议会致辞，随后提出三项立法建议，其草案被宣读。］}\par

在宣读之后，由我们至圣敬虔的君王将各条款交予最蒙天主眷爱的君士坦丁堡总主教阿纳托利（该城为新罗马），所有最敬爱天主的主教们喊道：愿皇帝与皇后陛下万岁，虔诚的基督徒。愿你们所事奉的基督保佑你们。这些事合乎信德。致司铎皇帝。你整顿了教会，击败了敌人，成为信德的导师。愿虔诚的皇后、爱基督者万岁。愿她这位正教徒万岁。愿天主保守你的王国。你除灭了异端者，持守了信德。愿仇恨远离你的帝国，你的王国永存！\par

我们至圣敬虔的君王对神圣主教会议说：为光荣圣殉道者欧斐弥并尊崇你们的圣德，我们钦定在此举行神圣信德会议的迦克墩城，得享大都市之荣衔，唯名义上赐此尊荣，而尼苛默底亚城原有的尊严仍予保留。\par

全体喊道，等等，等等。\par

\Emph{关于耶路撒冷与安提约基雅管辖权的法令}\par

% structure: p0068 source=h2 target=section reason=relative-heading-rank; base=h2

\section{第七次会议}

最显赫荣耀的法官们说：……根据最圣洁的安提约基雅主教马克西莫与最圣洁的耶路撒冷主教尤文纳勒双方的协议，如各自声明确认，此项安排藉我们之谕令及神圣主教会议之判决，应永为有效；即最圣洁的马克西莫主教，或更确切地说，最圣洁的安提约基雅教会，应管辖两个腓尼基与阿拉伯；而最圣洁的尤文纳勒主教，或更确切地说，其属下最圣洁的教会，应管辖三个巴勒斯坦地区；所有帝国诏令、文书及刑罚均予废止，遵照我们至圣敬虔君王的旨意。\par

巴勒里尼在其对圣良著作的注释中（米涅，《拉丁教父集》第五十五卷，第733页及以下）引用了本次会议记录的片段；若这些片段可信，则表明安提约基雅与耶路撒冷之权利问题曾在随后的一次会议（10月31日）上重新处理并以同样方式判定。这些片段通常被认为真实，并被曼西（第七卷，722 C.）收入其《公议会集》。\par

\Emph{关于厄弗所主教之法令}\par

% structure: p0072 source=h2 target=section reason=relative-heading-rank; base=h2

\section{第十二次会议}

最显赫的法官们说：鉴于敬爱天主的君士坦丁堡总主教阿纳托利与代表最敬爱天主的罗马总主教良之可敬主教帕斯加西努的提议——即因巴西安诺与斯德望二人行事均违反正典，故他们中谁亦不应管理，亦不得称为最圣洁的厄弗所教会之主教——并且全体神圣主教会议认为，这些祝圣乃违反正典而行，并赞同最可敬主教们的发言；因此，最可敬的巴西安诺与最可敬的斯德望将被免去厄弗所圣教会的职务；但得保留主教尊荣，并从上述最圣洁教会之收入中，每年各领取二百金币以作赡养与慰藉；另应按正典为最圣洁的教会祝圣另一位主教。\par

全体神圣主教会议喊道：这是公正的判决。这是敬虔的方案。这些事是公平的。\par

最可敬的巴西安诺主教说：请下令归还从我处被窃之物。\par

最光荣的法官说：\Quote{如果属于最可敬的\Person{巴西安}主教的任何个人东西被任何人（无论是最可敬的\Person{斯德望}主教还是任何其他人）从他那里取走，在经过司法证明后，应由取走或导致取走的人予以归还。}\par

\Emph{关于\Place{尼科美底亚}的决定}\par

% structure: p0078 source=h2 target=section reason=relative-heading-rank; base=h2

\section{第十三会议}

最光荣的法官说【在宣读御函结束后】：\Quote{这些神圣信件关于主教职位什么也没有说，但都提到属于都主教城市的荣誉。但神圣记忆的\Person{瓦伦提尼安}和\Person{瓦伦斯}的御函，当时授予\Place{尼斯}城市都主教权利，小心地规定不得从其他城市夺走任何东西。圣教父们的教规规定每省应有一个都主教。那么神圣主教会议在此事上的意愿是什么？}\par

圣主教会议高呼：\Quote{应遵守教规。让教规足够了。}\par

\Person{阿提克斯}，最可敬的\Place{伊庇鲁斯的旧尼科波利斯}主教，说：\Quote{教规如此规定：都主教应在各省有管辖权，并应设立该省所有主教。这就是教规的含义。如今\Place{尼科美底亚}主教，既然从一开始这就是一个都主教区，应祝圣该省所有主教。}\par

圣主教会议说：\Quote{这是我们都希望的，我们都祈求的，让这到处被遵守，这使我们所有人喜悦。}\par

\Person{若望}、\Person{君士坦丁}、\Person{帕特里克}（\Person{伯多禄}）以及\Person{若望}（他们中的一位）等最可敬的\Place{本都}教区主教们说：\Quote{教规承认更古老者为都主教。显然，最虔诚的\Place{尼科美底亚}主教享有祝圣权，并且既然法律（如阁下所见）仅以都主教之名尊崇\Place{尼斯}，从而使其主教仅在荣誉上高于该省其他主教。}\par

圣主教会议说：\Quote{他们按照教规教导，美丽地教导。我们都同意。}\par

\Emph{【\Person{埃提乌斯}，\Place{君士坦丁堡}总执事，随后提出诉求，以维护皇室城市宝座的权利。】}\par

最光荣的法官说：\Quote{最可敬的\Place{尼科美底亚}主教应拥有对\Place{比提尼亚}省教会的都主教权威，而\Place{尼斯}应仅享有都主教级别的荣誉，并按照\Place{尼科美底亚}省其他主教的榜样服从\Place{尼科美底亚}。因为这是神圣主教会议的意愿。}\par

% structure: p0087 source=h2 target=section reason=relative-heading-rank; base=h2

\section{第十五会议}

\Emph{神圣第四次大公会议（\Place{加采东}）的三十条教规。}\par

% structure: p0089 source=h3 target=subsection reason=relative-heading-rank; base=h2

\subsection{第一条教规}

\Quote{我们判定，至今为止每届会议制定的圣教父教规应继续有效。}\par

在\Place{加采东}大公会议召开之前，希腊教会中已被收集成册并连续编号的多次会议教规，正如我们所看到的，摆在了\Place{加采东}会议面前。然而，这些教规被收录的许多会议，例如\Place{新凯撒利亚}、\Place{安卡拉}、\Place{冈格拉}、\Place{安提约基雅}会议，显然不是大公会议，甚至在某种程度上权威存疑，比如341年的\Place{安提约基雅}会议。现在大公会议予以确认，以将它们提升为普遍且无条件有效的教会规则。皇帝\Person{犹斯提尼安}在其第131新律第一章中出色地指出：\Quote{我们尊崇前四次大公会议的教义法令如同圣经，但尊崇由它们发布或\Emph{认可}的教规则如同法律。}\par

似乎不可能确定这个列表具体包含哪些会议；但特鲁洛会议在其第二条教规中完全消除了这种歧义。\par

% structure: p0093 source=h3 target=subsection reason=relative-heading-rank; base=h2

\subsection{第二条教规}

\Quote{如果任何主教为了金钱祝圣，并出售不可出售的恩宠，为了金钱祝圣主教、乡村主教、司铎、执事或任何其他列入神职人员的人；或出于贪财提名管家、护法、管事或任何在教会名册上的人，让他被证实如此行为者丧失其品级；而被祝圣者不得因购买的祝圣或升迁获益；但应被革除他以金钱获得的尊位或职务。如果有人被发现进行如此可耻和非法交易，如果他是神职人员，应被从品级罢免；如果他在俗或隐修士，应受绝罚。}\par

然而，根据\Person{Van Espen}（他在此依靠\Person{Du Cange}），由\Quote{prosmonarios}或\Quote{mansionarius}，与由\Quote{oiconomos}一样，应理解为教会财产的\OriginalTerm{管家}{prosmonarios}。\par

% structure: p0096 source=h3 target=subsection reason=relative-heading-rank; base=h2

\subsection{第三条教规}

神圣主教会议【得知】某些列入神职人员的人因贪财成为他人财产的租赁者，订立世俗事务的契约，轻视事奉天主，并溜进世俗人的家中，因贪心而承担管理他们的财产。因此，伟大而神圣的会议决定，从此任何主教、神职人员或隐修士不得租赁财产，或从事商业，或投身世俗事务，除非他应法律召唤而无法避免地担任未成年人的监护人；或除非城市主教出于敬畏天主，委托他管理教会事务，或关照无助的孤儿、寡妇以及特别需要教会帮助的人。若此后有人违反这些规定，则应受教会处罚。\par

% structure: p0098 source=h3 target=subsection reason=relative-heading-rank; base=h2

\subsection{第四条教规}

凡真实真诚度隐修生活者，应被视为堪受尊敬；但因某些人以隐修生活为借口，在城市中随意往来，扰乱教会与政治事务，同时试图为自己建立隐修院；兹规定：任何人不得违背所在城市主教的意愿，在任何地方建造或建立隐修院或祈祷所；各城市及地区的隐修士应服从主教，安于宁静的生活，专务斋戒与祈祷，永久居留于其被安置的地方；不得干涉教会或世俗事务，也不得离开其隐修院参与此类事务；除非因紧急需要，经该城市主教指派。未经主人同意，不得收留任何奴隶进入隐修院成为隐修士。若有人违反此判决，我们规定应予以绝罚，以免天主的名受亵渎。城市主教必须妥善照管隐修院的需要。\par

赫费勒：与前一条教规一样，此条由马西安皇帝在第六次会议中提出，作为第一条，主教会议几乎逐字接受了皇帝提出的教规。此教规的起因似乎是尤提克倾向的隐修士，尤其是叙利亚的巴尔苏马斯，从第四次会议可见。他和他的隐修士作为尤提克派，拒绝服从他们怀疑有聂斯多略主义的主教的管辖。\par

% structure: p0101 source=h3 target=subsection reason=relative-heading-rank; base=h2

\subsection{第五条教规}

关于主教或神职人员从一城往另一城者，兹规定：圣教父们所颁布的教规仍应保持效力。\par

据赫费勒推测，主教们当时想到的是巴西安的案件，他在第十一次会议（10月29日）诉讼中声称自己被暴力逐出厄弗所教区。实际主教斯德望答辩说，巴西安并非被\Quote{祝圣}为该教区，而是侵入该教区并被公正驱逐。巴西安反驳说，他最初被祝圣为厄瓦萨教区的主教是暴力甚至野蛮的；他从未造访过厄瓦萨，因此他被任命到厄弗所并非转换教区。最终，主教会议以快刀斩乱麻的方式解决，下令选举一位新主教，巴西安和斯德望保留主教衔，并从教区收入中领取津贴（曼西，卷七，273页及\Emph{以下}）。\par

% structure: p0104 source=h3 target=subsection reason=relative-heading-rank; base=h2

\subsection{第六条教规}

无论是司铎、执事，还是任何神职等阶，均不得被祝圣为无固定职位者，除非被祝圣者被特别指派到城市或乡村的某一教堂、殉道纪念所或隐修院。若有人未经指派而被祝圣，神圣主教会议裁定，为祝圣者的侮辱，该祝圣无效，且该人不得在任何地方主持礼仪。\par

赫费勒。\par

显然，本条教规禁止所谓的\OriginalTerm{绝对祝圣}{absolute ordinations}，并要求每位神职人员在祝圣时必须被指定到一个具体的教堂。这里唯一认可的\OriginalTerm{资格}{titulus}是后来被称为\OriginalTerm{圣职资格}{titulus beneficii}的。作为此类资格的不同种类，我们在此看到：(a) 被任命到城市教堂；(b) 被任命到乡村教堂；(c) 被任命到殉道者小堂；(d) 被任命为隐修院的专职司铎。要正确理解最后一点，必须记住最早的隐修士绝非神职人员，但很快在每个较大的修院中，习惯上至少有一位隐修士被祝圣为司铎，以便在隐修院中提供礼仪服务。\par

布赖特。\par

通过\OriginalTerm{\GreekText{μαρτυρίῳ}}{martyry}一词，意指建在殉道者坟墓上的教堂或小堂。因此老底嘉主教会议禁止神职人员访问\Quote{异端者的殉道纪念所}（第九条教规）。尼撒的额我略谈到\Quote{神圣殉道者的殉道纪念所}（《著作》，卷二，212）；金口若望提到一个\Quote{殉道纪念所}，帕拉迪乌斯提到安提约基雅附近的\Quote{殉道纪念所}（《宗徒大事录讲道集》，三八，5；《对话录》，17页），帕拉迪乌斯提到君士坦丁堡的\Quote{圣若翰殉道纪念所}（《对话录》，25页）。参见苏格拉底，卷四，18、23，关于埃德萨的圣多默和罗马的圣伯多禄及圣保禄的\Quote{殉道纪念所}；以及卷六，6，关于加采东的圣欧斐米亚的\Quote{殉道纪念所}，主教会议实际在此举行。在可见见证的明确意义上，该词用于耶路撒冷的复活教堂（欧瑟伯，《君士坦丁传》，卷三，40，卷四，40；曼西，卷六，564；济利禄，《教理讲授》，十四，3），以及圣墓本身（《君士坦丁传》，卷三，28）。建在殉道者坟墓上的教堂在西方被称为\OriginalTerm{memoriæ martyrum}{殉道者纪念所}，见《非洲法典》第八十三条（比较奥古斯丁，《论对亡者的关怀》第六章）。\par

% structure: p0110 source=h3 target=subsection reason=relative-heading-rank; base=h2

\subsection{第七条教规}

我们规定：凡已登记在神职班中或已成为隐修士者，不得接受军事职务或任何世俗尊荣；若他们胆敢如此行而不悔改以致回转他们最初为爱天主所选择的道路，则应予以绝罚。\par

% structure: p0112 source=h3 target=subsection reason=relative-heading-rank; base=h2

\subsection{第八条教规}

贫民院、隐修院及殉道纪念所的神职人员，应依照圣教父的传统，留在各城市主教的权下；不得有人傲慢地摆脱自己主教的管辖；若有人以任何方式违反本条教规，不服从自己的主教，如为神职人员，应受教规处分；如为隐修士或平信徒，应受绝罚。\par

% structure: p0114 source=h3 target=subsection reason=relative-heading-rank; base=h2

\subsection{第九条教规}

若有圣职人员与另一圣职人员有争执，他不可离开自己的主教而转向世俗法庭；但应先将此事呈报给自己的主教，或在主教同意下，由双方选定某人处理此事。若有人违反这些法令，应受教会法规的处罚。若圣职人员控告自己的主教或其他主教，应由省区教务会议裁决。若主教或圣职人员与省区都主教有分歧，应诉诸该教区之监督，或帝国京城的君士坦丁堡宗座，并在那里受审。\par

% structure: p0116 source=h3 target=subsection reason=relative-heading-rank; base=h2

\subsection{第十条}

圣职人员不得同时登记在两座城市的教会中，即最初被祝圣的教会与另一因规模较大而因贪图虚誉迁往的教会。凡如此行者，应遣回原被祝圣的教会，只在那里服务。若有人此前已从一教会迁至另一教会，则不得干预原教会的事务，亦不得干预属于原教会的殉道纪念所、济贫院及客栈。若在此大公会议的法令之后，有人胆敢做任何现在禁止之事，会议规定他应被剥夺品级。\par

这里出现了一种新的机构，有多处实例。尤利安曾指示按基督徒模式建立异教客栈（\OriginalTerm{异教客栈}{\GreekText{ξενοδοχεῖα}}）（\WorkTitle{书信}四九）。该撒利亚的巴西略院既是\OriginalTerm{客栈}{\GreekText{ξενοδοχεῖον}}，也是\OriginalTerm{济贫院}{\GreekText{πτωχεῖον}}；它设有\OriginalTerm{旅客住宿处}{\GreekText{καταγώγια} \GreekText{τοῖς} \GreekText{ξένοις}}，也为行路者和因疾病需要援助者而设，巴西略区分了从事慈善服务的各类人员，包括护送旅客者（\OriginalTerm{护送者}{\GreekText{τοὺς} \GreekText{παραπέμποντας}}，\WorkTitle{书信}九四）。热罗尼莫写信给潘玛基乌斯说：\Quote{我听说你在罗马港建造了一座客栈；}并补充说他本人也为前往白冷的朝圣者建造了一处\Quote{旅舍}（\WorkTitle{书信}十六，11，14）。金口若望在君士坦丁堡提醒听众说\Quote{教会设有一处公共居所}，和\Quote{称为客栈}（\WorkTitle{宗徒大事录讲道集}四五，4）。他的朋友奥林比亚斯对\Quote{客栈}慷慨捐赠（\WorkTitle{拉乌西奥历史}144）。在尼特里亚隐修院教会附近有一座客栈（同上，7）。伊斯基里翁在迦克墩第三次会议上宣读的备忘录中，抱怨他的宗主教狄奥斯库若滥用一位慈善妇女在埃及遗赠给\OriginalTerm{客栈和济贫院}{\GreekText{τοῖς} \GreekText{ξενεῶσι} \GreekText{καὶ} \GreekText{πτωχείοις}}的款项，并说他本人被狄奥斯库若关在麻风病人的\Quote{客栈}里（曼西，卷六，1013，1017）。查士丁尼在\WorkTitle{法典}一，3，49中提及客栈，在\WorkTitle{新律}134，16中提及客栈管理人。大额我略命令客栈的账目应由主教审计（\WorkTitle{书信}四，27）。查理大帝规定修复破旧的\Quote{客栈}（803年敕令；佩尔茨，\WorkTitle{法律}一，110）；阿尔昆劝告他的学生、约克总主教埃恩巴尔德，考虑在约克教区何处可以建立\Quote{客栈，即医院}（\WorkTitle{书信}五十。）。\par

% structure: p0119 source=h3 target=subsection reason=relative-heading-rank; base=h2

\subsection{第十一条}

我们规定，贫困者和需要援助者经审查后，应从教会领取仅作和平信用的信函，而非推荐信；因为推荐信只能给予可疑之人。\par

% structure: p0121 source=h3 target=subsection reason=relative-heading-rank; base=h2

\subsection{第十二条}

据悉，某些人违反教会法律，求助于世俗权力，通过帝国敕令将一个省分为两个，以致一省有两个都主教；因此，神圣会议规定今后任何主教不得企图行此事，否则将被剥夺品级。但已通过帝国文书获得都会之名的城市及其负责主教，仅保留空衔，所有都主教权利仍归真正的都会。\par

% structure: p0123 source=h3 target=subsection reason=relative-heading-rank; base=h2

\subsection{第十三条}

未经自己主教推荐信的陌生或不知名圣职人员，绝对禁止在另一城市行使职务。\par

% structure: p0125 source=h3 target=subsection reason=relative-heading-rank; base=h2

\subsection{第十四条}

既然在某些省，读经员与歌咏员获准结婚，神圣会议规定他们中任何人不得娶异端女子为妻。但若有人已从这样的婚姻生育子女，且子女已在异端者中受洗，则须带他们加入天主教会之共融；若子女尚未受洗，今后不得在异端者中为他们施洗，亦不得将他们嫁给异端者、犹太人、或外教人，除非与正统子女结婚者承诺改信正统信仰。若任何人违犯此神圣会议的法令，应受教会法规的谴责。\par

% structure: p0127 source=h3 target=subsection reason=relative-heading-rank; base=h2

\subsection{第十五条}

\Emph{妇人}在四十岁之前不得接受覆手成为女执事，且须经过严格审查。若她在接受覆手后服务一段时间，竟蔑视天主的恩宠而结婚，则她与与她结合之男子均应受绝罚。\par

% structure: p0129 source=h3 target=subsection reason=relative-heading-rank; base=h2

\subsection{第十六条}

凡已献身于主天主的童贞女及隐修士，不得结婚；若发现有人这样做，应受绝罚。但我们规定，在任何地方，主教有权对他们施以宽免。\par

% structure: p0131 source=h3 target=subsection reason=relative-heading-rank; base=h2

\subsection{第十七条}

在各省内，偏远或乡村堂区应继续服从现今管辖它们的主教，尤其是主教们已和平而连续地治理它们达三十年之久者。但若在三十年内对此曾有或现在有争议，自认为受亏者可以合法地将案件提交省区会议。若有人被其都主教亏待，则此事务应由该教区的\OriginalTerm{督主教}{exarch}或由\Place{君士坦丁堡}的宗座裁决，如前所述。若有城市已经或将来因皇帝权柄而新建，则教会堂区的秩序应依照政治与市政的范例。\par

% structure: p0133 source=h3 target=subsection reason=relative-heading-rank; base=h2

\subsection{第十八条}

即使世俗法律也完全禁止共谋或结党的罪行，在天主教会中更应予以禁止。因此，若有任何人，无论是神职人员或隐修士，被发现共谋、结党，或策划反对他们的主教或同僚神职人员的阴谋，他们务必要从自身等级中被撤职。\par

% structure: p0135 source=h3 target=subsection reason=relative-heading-rank; base=h2

\subsection{第十九条}

既然我们听闻在各省区内，主教的正典会议并未举行，因此许多需要改革的教会事务被忽略；因此，按照圣教父们的教规，神圣会议决定：每省的主教应每年两次聚集在都主教所认可的地点，并处理所有已发生的事务。凡不参加会议，而是留在自己城市的主教，即使身体健康且无任何不可避免的必要事务，应接受兄弟般的劝诫。\par

% structure: p0137 source=h3 target=subsection reason=relative-heading-rank; base=h2

\subsection{第二十条}

正如我们已颁布的法令，在某一教会服务的神职人员不得被任命到另一城市的教会，而应坚守他们起初被认为堪当供职的教会；但那些因不得已而被迫离开本国、因而迁至另一教会者不在此限。若在此法令之后，有主教接受属于另一位主教的神职人员，则规定：被接受者与接受者均应被绝罚，直至该迁走的神职人员返回自己的教会为止。\par

% structure: p0139 source=h3 target=subsection reason=relative-heading-rank; base=h2

\subsection{第二十一条}

\Emph{神职人员}和教友若对主教或神职人员提出控告，作为控告人不应被轻率、不经审查地接纳，而应首先调查他们自己的品性。\par

% structure: p0141 source=h3 target=subsection reason=relative-heading-rank; base=h2

\subsection{第二十二条}

神职人员在其主教去世后，不得夺取属于主教的东西，正如古代教规也已禁止；凡如此行的人将面临从自身等级中被降级的危险。\par

% structure: p0143 source=h3 target=subsection reason=relative-heading-rank; base=h2

\subsection{第二十三条}

神圣会议听闻，某些神职人员和隐修士，未经他们自己主教的许可，有时甚至在被主教处以绝罚期间，前往帝都\Place{君士坦丁堡}，并在那里长期居留，制造骚乱，扰乱教会秩序，并将人们的家室搅得翻天覆地。因此，神圣会议决定：首先应由最神圣的君士坦丁堡教会的\OriginalTerm{辩护人}{Advocate}通知此类人离开帝都；若他们仍无耻地继续同样行为，则由同一辩护人强制驱逐他们，令其返回原处。\par

% structure: p0145 source=h3 target=subsection reason=relative-heading-rank; base=h2

\subsection{第二十四条}

\Emph{隐修院}，一旦经主教同意而被祝圣，应永远作为隐修院，其所属财产应予以保存，绝不能再成为世俗住所。凡允许此事发生者，应承担教会惩罚。\par

% structure: p0147 source=h3 target=subsection reason=relative-heading-rank; base=h2

\subsection{第二十五条}

\Emph{鉴于}某些都主教，如我们所听闻，忽视托付给他们的羊群，并拖延主教祝圣，神圣会议决定：主教祝圣应在三个月内进行，除非某次不可避免的必要情况要求延长延迟期限。若他未这样做，应承担教会惩罚，且该空缺教会的收入应由同一教会的总管保管。\par

% structure: p0149 source=h3 target=subsection reason=relative-heading-rank; base=h2

\subsection{第二十六条}

\Emph{鉴于}我们听闻在某些教会中，主教在没有总管的情况下管理教会事务，兹决议：凡有主教的教会，也应从其所属神职人员中选出一位总管，在其主教授权下管理教会事务；如此，教会行政管理便不乏见证，且教会财物不至被浪费，亦不致使司祭职蒙受羞辱；若主教不这样做，他应受神圣教规的制裁。\par

% structure: p0151 source=h3 target=subsection reason=relative-heading-rank; base=h2

\subsection{第二十七条}

神圣会议规定：凡以婚姻为借口强行掳掠妇女者，以及此类掠夺者的帮助者或教唆者，若是神职人员，应被降级；若是教友，应被绝罚。\par

% structure: p0153 source=h3 target=subsection reason=relative-heading-rank; base=h2

\subsection{第二十八条}

\Emph{遵从}圣教父们的所有决定，并承认刚才宣读的、由一百五十位蒙天主所爱的主教（曾在蒙福纪念的皇帝\Person{狄奥多西}时代于帝国京城\Place{君士坦丁堡}，即新\Place{罗马}，集会）所制定的法令，我们同样规定并颁布关于至圣\Place{君士坦丁堡}教会（即新\Place{罗马}）之特权的事宜。因为教父们正确地将特权授予旧\Place{罗马}的宗座，因为它是帝国之都。而一百五十位最虔诚的主教，受同样考虑驱使，将同等特权（\OriginalTerm{同等特权}{\GreekText{ἴσα} \GreekText{πρεσβεῖα}}）授予至圣的新\Place{罗马}宗座，公义地判断，那尊享王权与元老院、并与旧帝国\Place{罗马}享有同等特权的城市，在教会事务中也应如她一样被尊崇，并仅次于她；因此，在本都、亚细亚和色雷斯教区中，只有总主教以及上述教区中居住在蛮族地区的主教，应由前述至圣\Place{君士坦丁堡}至圣教会宗座祝圣；而上述教区的每位总主教，与其省区的主教们一同，祝圣其本省的主教，正如神圣法令所规定的；但如前所述，上述教区的总主教应在按照惯例进行适当选举并报告给他之后，由\Place{君士坦丁堡}总主教祝圣。\par

% structure: p0155 source=h3 target=subsection reason=relative-heading-rank; base=h2

\subsection{第二十九法令}

将一位主教贬至司铎之列是亵圣行为；但若他们因正当理由被免去主教职务，则也不应拥有司铎的地位；若他们未经任何指控而被免职，则应恢复其主教尊位。\par

\Person{Anatolius}，最可敬的\Place{君士坦丁堡}总主教说：\Quote{若那些据称被从主教尊位贬至司铎等级的人，确实因充足理由而被定罪，则他们显然不配司铎的尊荣。但若他们无端被降入较低等级，则若显得无过，便应重新获得主教职位的尊荣与司祭职。}\par

所有最可敬的主教们喊道：\par

\Quote{教父们的判决是正义的。我们众口一词。教父们作出了公义的裁决。让总主教们的裁决生效。}\par

最显赫尊荣的法官们说：\par

\Quote{愿神圣主教会议的旨意永立。}\par

\Emph{取自同一神圣主教会议第四次会议，涉及埃及主教事宜。}\par

最显赫尊荣的法官们及全体元老院说：\par

% structure: p0164 source=h3 target=subsection reason=relative-heading-rank; base=h2

\subsection{第三十法令}

既然埃及最虔诚的主教们目前推迟了对最神圣的\Person{良}总主教信函的签署，并非因为他们反对大公信仰，而是因为他们声称在埃及教区，未经其总主教的同意和命令不得行此类事，并请求在\Place{亚历山大}新任总主教祝圣之前予以谅解，我们认为合情合理且仁慈地应允他们，让他们在帝国京城中保留其正式身份，直至\Place{亚历山大}城总主教被祝圣为止。\par

最虔诚的主教\Person{Paschasinus}，\Place{罗马}宗座代表，说：\par

\Quote{若阁下建议并命令对他们施以任何宽待，则让他们提供保证，在\Place{亚历山大}城获得一位主教之前，他们不会离开此城。}\par

最显赫尊荣的法官们及全体元老院说：\par

\Quote{让最神圣的\Person{Paschasinus}的裁决得到确认。}\par

\Quote{因此，让他们（即埃及最虔诚的主教们）保持其正式身份，要么提供保证（若可能），要么受誓言约束。}\par

\Signature{Hefele}\par

此段与上一段一样，并非正式法令，而是帝国专员在第四次会议中提出的一项建议的口头重复，经代表\Person{Paschasinus}改进，并由主教会议批准。此外，这条所谓的法令在古代汇编中并未出现，很可能是以与前一条相同的方式和相同的原因被添加到二十八条法令之后的。\par

\Signature{Bright}\par

会议本可坦率地坚持先听后判的责任（见第二十九法令注释）；然而，在此场合下，主教们一个接一个地发表了严厉无情、专横的言论，其唯一借口在于：他们对449年8月8日\Person{Dioscorus}的暴行记忆犹新，而他们急于恐吓或压制的三位埃及请愿者，实际上在那悲惨的一天曾支持过\Person{Dioscorus}。人之常情难以忘怀此事；但这一结果却是\Place{加采东}大公会议历史上的一个污点。\par

% structure: p0175 source=h2 target=section reason=relative-heading-rank; base=h2

\section{第十六次会议}

\Person{Paschasinus}和\Person{Lucentius}，最可敬的主教们，宗座代表，说：\Quote{若阁下如此命令，我们有事呈报。}\par

最显赫的法官们说：\Quote{请讲。}最神圣的主教\Person{Paschasinus}，\Place{罗马}代表，说：\Quote{世界统治者们，关怀神圣大公信仰（其王国与荣耀由此增长），曾愿作出此界定，以使所有教会通过神圣平安保持合一。但他们的仁慈更关切地眷顾未来，以免在天主的主教们之间再次发生纷争、分裂或丑闻。然而，昨日在阁下及我等卑微者离开之后，据说作出了一些决定，我们认为有违教规及教会纪律。我们恳请阁下命人宣读这些决定，使全体弟兄得悉所行之事是否公正。}\par

最显赫的法官们说：\Quote{若我们离开后有任何事发生，请予宣读。}\par

且在宣读之前，\Person{君士坦丁堡教会}的总执事\Person{艾蒂乌斯}说：\Quote{有关信仰的事项业已妥为处理，这是确定的。但在主教会议中，历来如此：在最重要的各事决定之后，也应审查并整理其他必要事项。我们——我是指君士坦丁堡至圣教会——显然有需要处理的事。我们曾请求来自罗马的\OriginalTerm{主教大人们}{\GreekText{κυρίοις} \GreekText{τοις} \GreekText{ἐπισκοποις}}与我们一同处理这些事务，但他们拒绝了，声称未收到相关指示。我们将此事提交给阁下您，您便命至圣主教会议考虑此点。当阁下您退出后，因这事关乎共同利益，至圣主教们起立，祈求此事得以实行。他们当时在场，这一切不是在隐藏或秘密的方式下进行的，而是按部就班并符合教会法规。}\par

最荣耀的审判官们说：\Quote{请宣读会议记录。}\par

\EditorialNote{随后宣读了法规（编号第二十八号）以及全部192个签名，包括\Place{安提约基雅}、\Place{耶路撒冷}和\Place{赫拉克利亚}的主教，但不包括\Place{凯撒勒雅}的\Person{特拉修斯}（他后来同意了）。仅在一周前，已有350人签署了信仰界定。当最后一个名字宣读完毕时，爆发了如下辩论。（科卢姆 810。）}\par

最可敬的主教及宗座代表\Person{卢森提乌斯}说：\Quote{首先请阁下注意，此事是通过欺骗至圣主教们而实现的，使他们被迫签署了尚待书写的法规，即他们所提及的那些法规。}\EditorialNote{希腊文本略有不同（我遵循了拉丁文本，因批评家们认为它比我们现在拥有的希腊文本更纯正）：\Quote{阁下已经注意到，在主教们面前做了许多事，以免任何人被迫签署上述法规；由必要性而定。}}\par

最可敬的主教们喊道：\Quote{没有人被强迫。}\par

最可敬的主教及宗座代表\Person{卢森提乌斯}说：\Quote{显然，三百一十八位教父的法令被搁置了，只提及了一百五十位的法令，而那些法令在主教会议的法规中并无任何地位，且如他们所承认，是在八十年前制定的。因此，若他们在这些年中享有了这项特权，现在他们寻求什么？若他们从未使用过，为何要寻求？}\EditorialNote{希腊文本如下：\Quote{显然，当前的法令被附加到三百一十八位和其后一百五十位的法令上，这些法令并未被纳入主教会议的法规，他们声称这些是已界定的。因此，若他们在当时使用了这项益处，现在他们寻求什么根据法规他们未曾使用过的呢？}}\par

君士坦丁堡至圣教会的总执事\Person{艾蒂乌斯}说：\Quote{若他们在此事上收到了任何指令，请将之呈上。}\par

\Person{博尼法丘斯}，司铎及宗座代理，说：\Quote{至福的宗座教宗，除了其他事，给了我们这个命令。}他从文书上宣读：\Quote{圣教父的规定不可轻率地违反或削减。我个人的尊严须由你们以各种方式守护。若有人受自己城市权势的影响，企图进行侵占，应用适当的坚定态度抵制之。}\par

最荣耀的审判官们说：\Quote{让各方引用法规。}\par

---------------------------------\par

最可敬的主教及代表\Person{帕斯卡西努斯}宣读：三百一十八位圣教父的第六条法规。\Quote{罗马教会始终拥有首席权。因此，埃及应如此自持：\Place{亚历山大里亚}主教拥有对所有教会的权力，因为这同样是关于罗马主教的惯例。同样，在\Place{安提约基雅}及其他省，较大城市的教会应拥有首席权。\EditorialNote{希腊文本为\Quote{让首席权保留给教会。}这句话我不理解，除非它意味着为了教会的利益，必须维护\Place{安提约基雅}的首席权。但此种情绪人们会预期在拉丁文本中而非希腊文本中找到。}还有一件事十分清楚：若有人违背教省主教的意愿被祝圣为主教，此大公会议已规定，这样的人不应成为主教。然而，若所有他同僚的判断合理且符合法规，而有两三人因自身固执而反对，则多数派的投票应生效。因为一种习俗已经盛行，且有一古老传统：应尊崇\Place{耶路撒冷}主教，让他享有相应的尊荣，但他自己教省的权利必须保留。}\par

书记官\Person{康斯坦丁}从总执事\Person{艾蒂乌斯}递给他的书中宣读：三百一十八位圣教父的第六条法规。\Quote{让古代习俗继续有效，即埃及的习俗，以使\Place{亚历山大里亚}主教拥有对所有教会的管辖权，因为这在罗马也是惯例。同样，在\Place{安提约基雅}及其他省份，让教会保留\OriginalTerm{首席权}{\GreekText{πρεσβεῖα}}。因为这一点绝对清楚：若有人违背教省主教的意愿被祝圣为主教，此大公会议已规定，这样的人不应成为主教。然而，若所有人基于理性且符合法规的共同投票中，有两三人因自身固执而反对，则多数派的投票应成立。}\par

同一位书记官从同一部法典中宣读了第二次大公会议的决定。\Quote{这些主教因天主的恩宠从遥远的省份聚集在君士坦丁堡，制定了以下规定：……主教不可到其教区范围以外的教会去，也不可扰乱教会；但根据教规，亚历山大主教只应管理埃及的事务，东方主教只应治理东方教区，安提阿教会应按照尼西亚教规所应得的尊荣予以维护；亚洲主教只应照管亚洲教区，本都的主教只应照管本都的事务，色雷斯的主教只应照管色雷斯的事务。但主教不应未经邀请进入另一教区进行祝圣或任何其他教会职务。在遵守上述关于教区的教规的情况下，显然每个省的教务会议将按照尼西亚的决定管理该省的事务。天主在外邦中的教会必须按照从教父时代以来一直盛行的习惯来管理。然而，君士坦丁堡的主教应拥有仅次于罗马主教的尊荣特权，因为君士坦丁堡是新罗马。}\par

---------------------------------\par

最荣耀的审判官说：让刚刚签署了所读卷册的最神圣的亚洲和本都主教表明，他们是凭自己的判断签署，还是因任何必要而被迫签署。当他们走到中间时，可敬的基齐库斯主教狄奥吉尼斯说：我呼求天主作证，我是凭自己的判断签署的。\EditorialNote{如此等等，一个接一个。}\par

其余的人喊道：我们是自愿签署的。\par

最荣耀的审判官说：既然显然每一位主教的签署都不是出于任何必要，而是出于自己的意愿，那么让尚未签署的最神圣的主教们发言。\par

安卡拉主教尤西比乌斯说：我只为自己发言。\par

\EditorialNote{他的发言是对自己为冈格拉祝圣主教一事所作的个人解释。}\par

最荣耀的审判官说：根据已经完成并提出的事实，我们认为，按照教规，万有之首的首席权（\OriginalTerm{首席权}{\GreekText{πρὸ} \GreekText{πάντων} \GreekText{τὰ} \GreekText{πρωτεῖα}}）和独特的尊荣（\OriginalTerm{独特的尊荣}{\GreekText{τὴν} \GreekText{ἐξαίρετον} \GreekText{τιμὴν}}）应归于旧罗马最受天主爱戴的总主教，但皇城君士坦丁堡（即新罗马）最可敬的总主教应享有同样的首席尊荣，并有权祝圣亚洲、本都和色雷斯教区的都主教，方式如下：由各都城的圣职人员、有产者（\OriginalTerm{有产者}{\GreekText{κτητόρων}}）和最杰出的人士，此外还有该省所有最可敬的主教或其多数，进行选举；被选者应为上述人员认为配得都主教职位的人；然后由所有选举他的人将其呈交给最神圣的皇城君士坦丁堡的总主教，以询问他（即君士坦丁堡宗主教）是否愿意在那里祝圣他，或委托他在获得投票选举的地方所在省祝圣他。普通城市的最可敬的主教应由该省所有最可敬的主教或多数祝圣，都主教根据教父制定的既定教规行使权力，并且对于此类祝圣，不向最神圣的皇城君士坦丁堡的总主教通报。如此，事情在我们看来就是如此。愿神圣大公会议惠赐教导其对本案的看法。\par

最可敬的主教们喊道：这是一个公正的判决。我们都这样说。这些事令我们都喜悦。这是一个公正的决定。确立所提议的法令形式。这是一个公正的表决。一切都按应当的方式颁布。我们恳求你让我们离开。凭着皇帝的安全，让我们离开。我们都会保持这个意见，我们都这样说。\par

Lucentius主教说：宗座下令所有事情都应在我们面前进行\EditorialNote{这句话在拉丁文中读作：宗座不应在我们面前受屈辱。我不知道为什么布莱特教规在他的第28条教规注释中采纳了这个读法}；因此，凡昨天在我们缺席期间所做有损于教规的事，我们恳请殿下下令撤消。否则，请将我们的反对意见记录在案，并请让我们清楚地知道\EditorialNote{拉丁文：我们得以知道}我们应当向那位最宗徒性且是全体教会统治者的主教报告什么，以便他能够就对其教座所施的侮辱和对教规的蔑视采取行动。\par

\EditorialNote{塞巴斯特最可敬的主教若望说：我们都会保持阁下所表达的意见。}\par

最荣耀的审判官说：整个大公会议已经批准了我们所提出的。\par

\SourceRecord{Translated by Henry Percival.；Nicene and Post-Nicene Fathers, Second Series；Vol. 14.；Edited by Philip Schaff and Henry Wace.；Buffalo, NY: Christian Literature Publishing Co.,；1900.；<http://www.newadvent.org/fathers/3811.htm>.}
